\documentclass[border=10pt]{standalone}
\usepackage{circuitikz}
\usepackage{xcolor}
\definecolor{myorange}{HTML}{E69F00}
\definecolor{mypurple}{HTML}{7B61FF}
\usepackage{amsmath}

% Reuse custom CircuitikZ styles/symbols (from jk_tikz_elec_circuit (3).tex)
\makeatletter
% Do not add spurious spaces

\tikzset{jk/.style=
      {%
            % generic tikz-level settings
            cute inductors,%
            %circutikz-level settings
            \circuitikzbasekey/.cd,%
            bipoles/thickness=1.5,
            bipole voltage style/.style={color=red},
            bipole current style/.style={color=blue},
            % Capacitors
            capacitors/thickness=1.2,
            capacitors/width/.try=0.085,
            capacitors/height/.try=0.3,
            % Inductors (and transformers)
            inductors/scale=0.9,
            inductors/thickness=1.0,
            inductors/coils=4,
            inductors/width/.try=0.6,
            % Ground symbols
            grounds/scale=0.6,
            grounds/thickness=1.0,
            % Nodes width (connectors, etc)
            nodes width=.04,
            % Arrows size
            current arrow scale=16,
            % Transmission line
            bipoles/tline/height=0.5,
      },
}%

%%%%%%%%%%%%%%%%%%%%%%%%%%%%%%%%%%%%%%%%%%%%%%%%%%%%%%%%%%%%%%%%%%%%%%%%%%%
%% JJ (junction)
\ctikzset{bipoles/jj/height/.initial=.23}
\ctikzset{bipoles/jj/width/.initial=.23}
\ctikzset{bipoles/jj/thickness/.initial=1}

\pgfcircdeclarebipolescaled{inductors}
      {} % extra anchors
      {\ctikzvalof{bipoles/jj/height}} % depth under path line
      {jj} % name
      {\ctikzvalof{bipoles/jj/height}} % height above path line
      {\ctikzvalof{bipoles/jj/width}}  % width
      { % code
            % cross
            \pgfsetlinewidth{\ctikzvalof{bipoles/jj/thickness}\pgfstartlinewidth}
            \pgfpathmoveto{\pgfpoint{\pgf@circ@res@left}{\pgf@circ@res@up}}
            \pgfpathlineto{\pgfpoint{\pgf@circ@res@right}{\pgf@circ@res@down}}

            \pgfpathmoveto{\pgfpoint{\pgf@circ@res@right}{\pgf@circ@res@up}}
            \pgfpathlineto{\pgfpoint{\pgf@circ@res@left}{\pgf@circ@res@down}}

            \pgfpathmoveto{\pgfpoint{\pgf@circ@res@left}{\pgf@circ@res@up}}
            \pgfpathlineto{\pgfpoint{\pgf@circ@res@right}{\pgf@circ@res@up}}
            \pgfpathlineto{\pgfpoint{\pgf@circ@res@right}{\pgf@circ@res@down}}
            \pgfpathlineto{\pgfpoint{\pgf@circ@res@left}{\pgf@circ@res@down}}
            \pgfpathlineto{\pgfpoint{\pgf@circ@res@left}{\pgf@circ@res@up}}
            \pgfusepath{draw}
      }

\pgfcirc@activate@bipole@simple{l}{jj}
\makeatother

\begin{document}
\begin{circuitikz}[jk, scale=0.85, transform shape]


% Left branch - LC resonator (cavity) - C and L in parallel (myorange)
\draw[color=myorange] (0,0) node[sground]{} to[short] (0,1.2);
\draw[color=myorange] (0,1.2) to[short] (-0.35,1.2) to[C] (-0.35,2.8) to[short] (0,2.8);
\draw[color=myorange] (0,1.2) to[short] (0.35,1.2) to[L] (0.35,2.8) to[short] (0,2.8);
\draw[color=myorange] (0,2.8) to[short] (0,4);


% --- Split-color capacitor between (0,4) and (2.2,4) ---
% Use the same width/height values you set in tikzset:
\def\Cgapmain{0.11}  % == capacitors/width
\def\Chmain{0.22}      % == capacitors/height
\def\Cwmain{0.5pt}
% Center of the capacitor
\def\xcmain{1.1}
\def\ycmain{4}
% Left lead (orange)
\draw[color=myorange] (0,4) -- (\xcmain-0.5*\Cgapmain,4);
% Left plate (orange)
\draw[color=myorange,line width=\Cwmain] (\xcmain-0.5*\Cgapmain,\ycmain-\Chmain) -- (\xcmain-0.5*\Cgapmain,\ycmain+\Chmain);
% Right plate (purple)
\draw[color=mypurple,line width=\Cwmain] (\xcmain+0.5*\Cgapmain,\ycmain-\Chmain) -- (\xcmain+0.5*\Cgapmain,\ycmain+\Chmain);
% Right lead (purple)
\draw[color=mypurple] (\xcmain+0.5*\Cgapmain,4) -- (2.2,4);


% Transmon (SQUID loop with flux) - same size and position as LC loop (mypurple)
\draw[color=mypurple] (2.2,4) to[short] (2.2,2.8);
\draw[color=mypurple] (2.2,2.8) to[short] (1.85,2.8) to[jj] (1.85,1.2) to[short] (2.2,1.2);
\draw[color=mypurple] (2.2,2.8) to[short] (2.55,2.8) to[jj] (2.55,1.2) to[short] (2.2,1.2);
\draw[color=mypurple] (2.2,1.2) to[short] (2.2,0);
\draw[color=mypurple] (2.2,0) node[sground]{};

% Flux label
\node at (2.2,2) {$\large\Phi$};


% Readout line (top right, mypurple)
\draw[color=mypurple] (2.2,4) to[short] (3.,4);

% --- Split-color capacitor that matches CircuitikZ settings ---
% Use the same width/height values you set in tikzset:
\def\Cgap{0.11}  % == capacitors/width
\def\Ch{0.22}      % == capacitors/height
% Choose a line width consistent with your circuitikz style.
% (If you want it identical to other caps, keep it modest; avoid 'very thick'.)
\def\Cw{0.5pt}


% capacitor centered at x=3.3, y=4.0
\draw (3.,4) -- (3.3-0.5*\Cgap,4);                         % left lead
\draw[line width=\Cw] (3.3-0.5*\Cgap,4-\Ch) -- (3.3-0.5*\Cgap,4+\Ch); % left plate (black)
\draw[color=mypurple,line width=\Cw] (3.3+0.5*\Cgap,4-\Ch) -- (3.3+0.5*\Cgap,4+\Ch); % right plate (mypurple)
\draw[color=mypurple] (3.3+0.5*\Cgap,4) -- (3.8,4);                   % right lead (mypurple)

% Voltage source (drive)
\draw[blue] (3.8,4) to[vsourcesin, l=$V_g$] (4.8,4);

% Output (grounded)
\draw[blue] (4.8,4) to[short] (5.4,4);
\draw[blue] (5.4,4) to[short] (5.4,3.2) node[sground]{};

% Flux noise label
\node[gray] at (3.2,2) {$\delta\Phi(t)$};

\end{circuitikz}
\end{document}
